\subsection{Three Dimensional Coordinate Systems}

\content{
}{% presentation
\usetikzlibrary{3d}
\hfill\begin{tikzpicture}[x={(-135:.9cm)}, y={(0:1cm)}, z={(90:1cm)}]
\draw[->] (0,0,0) -- (3.4,0,0);
\draw[->] (0,0,0) -- (0,3.4,0);
\draw[->] (0,0,0) -- (0,0,3.4);
\foreach \x/\xlabel in {.5/1,1/2,1.5/3,2/4,2.5/5,3/6} {
    \draw(\x,-.1,0)--(\x,.1,0);
    \node at (\x,0,0) [left=3pt] {\xlabel};
}
\foreach \y/\ylabel in {.5/1,1/2,1.5/3,2/4,2.5/5,3/6} {
    \draw(-.1,\y,0)--(.1,\y,0);
    \node at (0,\y,0) [above=3pt] {\ylabel};
}
\foreach \z/\zlabel in {.5/1,1/2,1.5/3,2/4,2.5/5,3/6} {
    \draw(0,-.1,\z)--(0,.1,\z);
    \node at (0,0,\z) [left=3pt] {\zlabel};
}
\vspace{.5in}
\end{tikzpicture}
}{% nopresentation
}{% comments
We describe 3-dimensional space by picking an origin, and three perpendicular
axis, namely $x$, $y$, and $z$-axis.  these must follow the right hand rule. 

For any point in 3-dimentional space, these three axis give us a way to describe
where that point is, namely a magnitude in each of the three directions of the
axes. We call this the coordinate of the point. 

\textbf{Give an example of plotting a point in 3-space.}
}

\content{
\subsubsection{Surfaces and Solids}
\begin{example} Describe the surfaces defined by the following equations. 
\begin{itemize}
    \item $z=3$
    \item $y=5$
\end{itemize}
\end{example}
}{% presentation
\usetikzlibrary{3d}
\begin{tikzpicture}[x={(-135:.9cm)}, y={(0:1cm)}, z={(90:1cm)}]
\draw[->] (0,0,0) -- (3.4,0,0);
\draw[->] (0,0,0) -- (0,3.4,0);
\draw[->] (0,0,0) -- (0,0,3.4);
\foreach \x/\xlabel in {.5/1,1/2,1.5/3,2/4,2.5/5,3/6} {
    \draw(\x,-.1,0)--(\x,.1,0);
    \node at (\x,0,0) [left=3pt] {\xlabel};
}
\foreach \y/\ylabel in {.5/1,1/2,1.5/3,2/4,2.5/5,3/6} {
    \draw(-.1,\y,0)--(.1,\y,0);
    \node at (0,\y,0) [above=3pt] {\ylabel};
}
\foreach \z/\zlabel in {.5/1,1/2,1.5/3,2/4,2.5/5,3/6} {
    \draw(0,-.1,\z)--(0,.1,\z);
    \node at (0,0,\z) [left=3pt] {\zlabel};
}
\vspace{.5in}
\end{tikzpicture}
% second diagram for the other problem
\usetikzlibrary{3d}
\hspace{.5in}\begin{tikzpicture}[x={(-135:.9cm)}, y={(0:1cm)}, z={(90:1cm)}]
\draw[->] (0,0,0) -- (3.4,0,0);
\draw[->] (0,0,0) -- (0,3.4,0);
\draw[->] (0,0,0) -- (0,0,3.4);
\foreach \x/\xlabel in {.5/1,1/2,1.5/3,2/4,2.5/5,3/6} {
    \draw(\x,-.1,0)--(\x,.1,0);
    \node at (\x,0,0) [left=3pt] {\xlabel};
}
\foreach \y/\ylabel in {.5/1,1/2,1.5/3,2/4,2.5/5,3/6} {
    \draw(-.1,\y,0)--(.1,\y,0);
    \node at (0,\y,0) [above=3pt] {\ylabel};
}
\foreach \z/\zlabel in {.5/1,1/2,1.5/3,2/4,2.5/5,3/6} {
    \draw(0,-.1,\z)--(0,.1,\z);
    \node at (0,0,\z) [left=3pt] {\zlabel};
}
\vspace{.5in}
\end{tikzpicture}
}{% nopresentation
}{% comments
Mention that in 2-dimensional geometry, graphs of equations involving $x$ and
$y$ are curves in the plane, while here, graphs of equations involving $x$, $y$,
and $z$ are surfaces. 
}

\content{
\begin{example} Describe the set of points which satisfy the two equations
$$ x^2+y^2=1 \quad \text{and}\quad z=3. $$
\end{example}
}{% presentation
\usetikzlibrary{3d}
\hfill\begin{tikzpicture}[x={(-135:.9cm)}, y={(0:1cm)}, z={(90:1cm)}]
\draw[->] (0,0,0) -- (3.4,0,0);
\draw[->] (0,0,0) -- (0,3.4,0);
\draw[->] (0,0,0) -- (0,0,3.4);
\foreach \x/\xlabel in {.5/1,1/2,1.5/3,2/4,2.5/5,3/6} {
    \draw(\x,-.1,0)--(\x,.1,0);
    \node at (\x,0,0) [left=3pt] {\xlabel};
}
\foreach \y/\ylabel in {.5/1,1/2,1.5/3,2/4,2.5/5,3/6} {
    \draw(-.1,\y,0)--(.1,\y,0);
    \node at (0,\y,0) [above=3pt] {\ylabel};
}
\foreach \z/\zlabel in {.5/1,1/2,1.5/3,2/4,2.5/5,3/6} {
    \draw(0,-.1,\z)--(0,.1,\z);
    \node at (0,0,\z) [left=3pt] {\zlabel};
}
\vspace{.5in}
\end{tikzpicture}
}{% nopresentation
}{% comments
}

\content{
\begin{definition} (Distance Formula) Given two points $P_1(x_1,y_1,z_1)$ and 
$P_2(x_2,y_2,z_2)$, the distance between them is 
$$ d(P_1,P_2) = \sqrt{(x_1-x_2)^2+ (y_1-y_2)^2+(z_1-z_2)^2}. $$
\end{definition}

\begin{example} Find the distance between the two points $P_1(2,-1,7)$ and
$P_2(1,-3,5)$.
\end{example}
}{% presentation
\vspace{1in}
}{% nopresentation
}{% comments
}

\content{
\begin{definition} (Equation of a Sphere) An equation of the sphere with center
$P(a,b,c)$ and radius $r\geq 0$ is given by 
$$ (x-a)^2 + (y-b)^2 + (z-c)^2  = r^2 . $$
\end{definition}
}{% presentation
}{% nopresentation
}{% comments
}

\content{
\begin{example} Find the equation of a sphere with center $(1,1,1)$ and passes
through $(2,-1,2)$.
\end{example}
}{% presentation
\vspace{1in}
}{% nopresentation
}{% comments
}

\content{
\begin{example} Show that the equation 
$$ x^2+y^2+z^2 + 4x-6y = 0 $$
is the equation of a sphere, and find it's center and radius. 
\end{example}
}{% presentation
\vspace{2in}
\newpage % end the page here, possibly for a video split. 
}{% nopresentation
}{% comments
}

\subsection{Vectors}


\content{
}{% presentation
\vspace{3in}
}{% nopresentation
}{% comments
Give a geometric description of vectors. Including scale multiples, and 
the parallelogram rule for adding vectors. Demonstrate vector subtraction
geometrically. Do several examples of combining these operations.
}

\content{
It will be convenient to use vectors along side coordinate systems. If we use
our $x$, $y$, and $z$-axes from before, and we place a vector $\vec{v}$ at the
origin, the terminal end of the vector has coordinates which we will use to
describe the vector itself. 
}{% presentation
\vspace{2in}
}{% nopresentation
}{% comments
Give a few examples of vectors (not necessarily starting at the origin), and
their coordinates.  Introduce $\langle\cdot \rangle $ notation for vectors. 

Give an example of finding a vector between two points in space. 
}

\content{
Each vector has a length associated to it, this can be calculated by the
following formula. 
\begin{definition} The length of a vector $\vec{v} = \langle a_1, a_2\rangle$ is
obtained by the formula 
$$ \sqrt{a_1^2+a_2^2}, $$
and is denoted $|\vec{v}|$. Similarly, for a 3-dimensional vector 
$\vec{u} = \langle a_1, a_2, a_3\rangle$, the length is given by 
$$ |\vec{u}| = \sqrt{a_1^2+a_2^2+a_3^2}. $$
\end{definition}
}{% presentation
}{% nopresentation
}{% comments
}

\content{
When using vectors with a coordinate system, we have a convenient way of
describing the vector operations of addition and scalar multiplication: 
\begin{definition}
If $\vec{v} = \langle a_1,a_2,a_3\rangle$, and $\vec{u} = \langle
b_1,b_2,b_3\rangle$, then 
$$ \vec{v} +\vec{u} = \langle a_1+b_1, a_2+b_2,a_3+b_3\rangle, $$ 
and if $c$ is a scalar, then 
$$ c\vec{v} = \langle ca_1, ca_2,ca_3\rangle. $$
\end{definition}
}{% presentation
}{% nopresentation
}{% comments
}

\content{
\begin{theorem}
if $\vec{a}$, $\vec{b}$, and $\vec{c}$ are vectors, and $c$ and $d$ are scalars,
then the following properties hold: 
\begin{enumerate}
    \item $\vec{a}+ \vec{b} = \vec{b} + \vec{a}$. 
    \item $\vec{a} + (\vec{b}+\vec{c}) = (\vec{a}+\vec{b})+\vec{c}$
    \item $\vec{a} + \vec{0} = \vec{a}$
    \item $\vec{a} + (-\vec{a}) = \vec{0}$
    \item $c(\vec{a}+ \vec{b}) = c\vec{a}+c\vec{b}$
    \item $(c+d)\vec{a} = c\vec{a} + d\vec{a}$
    \item $(cd)\vec{a} = c(d\vec{a})$
    \item $1\vec{a} = \vec{a}$.
\end{enumerate}
\end{theorem}
}{% presentation
}{% nopresentation
}{% comments
}


\content{
\begin{example} If $\vec{a} = \langle 4,0,3\rangle$ and $\vec{b} = \langle
-2,1,5\rangle$, find $a+b$, $a-2b$, $3a$, and $|a|$.
\end{example}
}{% presentation
\vspace{2in}
}{% nopresentation
}{% comments
}


\content{
For 3-dimensional space, there are 3 vectors which play a special role, namely
$$ \vec{i} = \langle 1,0,0\rangle,\quad \vec{j} = \langle 0,1,0\rangle,\quad
\vec{k} = \langle 0,0,1\rangle. $$
These are called the standard basis vectors, and we will often express vectors
in terms of $\vec{i}$, $\vec{j}$, and $\vec{k}$.  For example, we have that 
if $\vec{a} = \langle 3,-2,1\rangle$, then 
$$\vec{a} = \langle 3, -2, 1\rangle  = 3\vec{i} -2\vec{j} + \vec{k}. $$
}{% presentation
}{% nopresentation
}{% comments
}

\content{
\begin{example} For the vectors $\vec{a} = \vec{j} - \vec{k}$ and 
$\vec{b} = 3\vec{i} + 2\vec{j}$, find $\vec{a}+ 2\vec{b}$.
\end{example}
}{% presentation
\vspace{1in}
}{% nopresentation
}{% comments
}

\content{
\begin{definition} A unit vector is a vector whose length is 1.
\end{definition}
The vectors
$\vec{i}$, $\vec{j}$ and $\vec{k}$ are all unit vectors. If we have a vector
$\vec{a}$, and we wish to find a unit vector in the same direction as $\vec{a}$,
we compute 
$$ \vec{u} = \frac{\vec{a}}{|\vec{a}|}. $$
\begin{example} Find a unit vector in the direction of $\langle 2,3,-1\rangle$.
\end{example}
}{% presentation
\vspace{1in}
}{% nopresentation
}{% comments
}
